\sectionkicker{Source-backed technical notes}
\subsection*{知識を重みに閉じ込めない――REALMからRAGへ: Technical Notes}
本欄は記事本文で圧縮した一次資料上の情報を、比較・再検証しやすい形で整理したものである。時系列、確認済みの主張、留意点、一次資料URLを掲載する。
\smallskip
\noindent{\small\textit{Technical Notesでは、一次資料から確認した公開・提供・研究上の事実を記載する。数値・能力評価や提供元・著者による比較表現は、独立再現済みの結果ではなく帰属claimとして扱う。}}\par
\medskip
\subsection*{Theme at a glance}
\addcontentsline{toc}{subsection}{Theme at a glance}
\begin{center}
\footnotesize
\begin{tabularx}{\linewidth}{@{}p{0.20\linewidth}p{0.10\linewidth}p{0.14\linewidth}X@{}}
\toprule
資料 & 位置づけ & 種別 & 時系列 \\
\midrule
Retrieval-Augmented Generation (RAG) & 主要資料 & 論文 & 2020-05-22, 2020-09-28 \\
REALM & 補足資料 & 論文 & 2020-02-10 \\
\bottomrule
\end{tabularx}
\end{center}
\normalsize

\begin{technicalnote}{Retrieval-Augmented Generation (RAG)}{主要資料}
\begingroup
% reader-facing Technical Notes break policy
\widowpenalty=10000
\clubpenalty=10000
\displaywidowpenalty=10000
\begin{tabularx}{\linewidth}{@{}>{\bfseries}p{0.22\linewidth}X@{}}
組織 & Facebook AI Research / authors \\
種別 & 論文 \\
時系列 & 2020-05-22, 2020-09-28 \\
\end{tabularx}
\smallskip
\begin{minipage}{\linewidth}
% reader-facing Technical Notes event-only fact/source tail group
{\bfseries 一次資料から整理したtechnical points}
\begin{itemize}[leftmargin=1.5em,itemsep=0.35em]
\item \textbf{一次情報で確認できる事実}: Facebook AI Research / authorsの選定済み一次資料では、RAGはpretrained seq2seq modelのparametric memoryと、pretrained neural retrieverでアクセスするWikipediaのdense vector indexというnon-parametric memoryを組み合わせる。 論文はretrieved passagesをgenerated sequence全体で共有する形式と、tokenごとに異なるpassageを利用できる形式の2構成を比較する。 Metaの2020-09-28一次資料はDPRとBARTを組み合わせるend-to-end differentiable構成を説明し、RAGをHugging Face Transformers componentとして公開したことをavailability boundaryとして記録する。
\end{itemize}
\begin{samepage}
{\bfseries 一次資料}
\begingroup\sloppy
\Urlmuskip=0mu plus 2mu\relax
\begin{itemize}[leftmargin=1.5em,itemsep=0.25em]
\item {\scriptsize\url{http://arxiv.org/abs/2005.11401v4}}
\item {\scriptsize\url{https://ai.meta.com/blog/retrieval-augmented-generation-streamlining-the-creation-of-intelligent-natural-language-processing-models/}}
\end{itemize}
\endgroup
\end{samepage}
\end{minipage}
\endgroup
\end{technicalnote}
\medskip

\begin{technicalnote}{REALM}{補足資料}
\begingroup
% reader-facing Technical Notes break policy
\widowpenalty=10000
\clubpenalty=10000
\displaywidowpenalty=10000
\begin{tabularx}{\linewidth}{@{}>{\bfseries}p{0.22\linewidth}X@{}}
組織 & Google / authors \\
種別 & 論文 \\
時系列 & 2020-02-10 \\
\end{tabularx}
\smallskip
\begin{minipage}{\linewidth}
% reader-facing Technical Notes event-only fact/source tail group
{\bfseries 一次資料から整理したtechnical points}
\begin{itemize}[leftmargin=1.5em,itemsep=0.35em]
\item \textbf{一次情報で確認できる事実}: Google / authorsの選定済み一次資料では、REALMはlanguage-model pretrainingへlatent knowledge retrieverを組み込み、Wikipedia等のlarge corpusからdocumentをretrieveしてpretraining・fine-tuning・inferenceの各段階でattentionする。 retriever自体をmasked language modelingのlearning signalでunsupervised pretrainし、millions of documentsを考慮するretrieval stepを通してbackpropagationする設計を示した。 Open-domain QA上の性能差は著者実験として扱い、外部知識をparameterから分離するmodularityとinterpretabilityというarchitecture上の性質と区別する。
\end{itemize}
\begin{samepage}
{\bfseries 一次資料}
\begingroup\sloppy
\Urlmuskip=0mu plus 2mu\relax
\begin{itemize}[leftmargin=1.5em,itemsep=0.25em]
\item {\scriptsize\url{http://arxiv.org/abs/2002.08909v1}}
\end{itemize}
\endgroup
\end{samepage}
\end{minipage}
\endgroup
\end{technicalnote}
\medskip
